\documentclass[12pt,a4paper]{article}

% ── Encoding & Language ─────────────────────────────────────────────────────
\usepackage[T5]{fontenc}
\usepackage[utf8]{inputenc}
\usepackage[vietnamese]{babel}

% ── Page layout ─────────────────────────────────────────────────────────────
\usepackage[top=2.5cm,bottom=2.5cm,left=3cm,right=2.5cm]{geometry}
\usepackage{parskip}
\setlength{\parindent}{0pt}

% ── Typography ───────────────────────────────────────────────────────────────
\usepackage{microtype}
\usepackage{lmodern}

% ── Colors ───────────────────────────────────────────────────────────────────
\usepackage[dvipsnames,table]{xcolor}
\definecolor{hust}{RGB}{0,82,147}
\definecolor{rowgray}{RGB}{245,245,245}
\definecolor{codered}{RGB}{192,0,0}
\definecolor{codeblue}{RGB}{0,0,180}

% ── Headers & footers ────────────────────────────────────────────────────────
\usepackage{fancyhdr}
\pagestyle{fancy}
\fancyhf{}
\fancyhead[L]{\small\color{hust}Báo Cáo Bài Tập Lớn — ML \& Khai Phá Dữ Liệu}
\fancyhead[R]{\small\color{hust}HUST 2026}
\fancyfoot[C]{\thepage}
\renewcommand{\headrulewidth}{0.4pt}

% ── Section styling ──────────────────────────────────────────────────────────
\usepackage{titlesec}
\titleformat{\section}{\large\bfseries\color{hust}}{}{0em}{\thesection\quad}[\titlerule]
\titleformat{\subsection}{\normalsize\bfseries\color{hust!80!black}}{\thesubsection}{1em}{}
\titleformat{\subsubsection}{\normalsize\bfseries}{\thesubsubsection}{1em}{}

% ── Tables ────────────────────────────────────────────────────────────────────
\usepackage{booktabs}
\usepackage{tabularx}
\usepackage{longtable}
\usepackage{array}
\usepackage{multirow}
\usepackage{adjustbox}
\usepackage{ragged2e}
\usepackage{seqsplit}
\renewcommand{\arraystretch}{1.25}
\setlength{\tabcolsep}{4pt}
\newcolumntype{Y}{>{\RaggedRight\arraybackslash}X}
\newcolumntype{P}[1]{>{\RaggedRight\arraybackslash}p{#1}}
\newenvironment{fittabular}[2][c]{%
  \begin{adjustbox}{max width=\textwidth}%
  \begin{tabular}[#1]{#2}%
}{%
  \end{tabular}%
  \end{adjustbox}%
}

% ── Code blocks ──────────────────────────────────────────────────────────────
\usepackage{listings}
\lstset{
  basicstyle=\ttfamily\small,
  backgroundcolor=\color{rowgray},
  frame=single,
  framesep=4pt,
  xleftmargin=8pt,
  breaklines=true,
  breakatwhitespace=false,
  keepspaces=true,
}

% ── Lists ─────────────────────────────────────────────────────────────────────
\usepackage{enumitem}
\setlist[itemize]{topsep=4pt,itemsep=2pt}

% ── Math ─────────────────────────────────────────────────────────────────────
\usepackage{amsmath}
\usepackage{amssymb}

% ── Figures ──────────────────────────────────────────────────────────────────
\usepackage{graphicx}
\usepackage{caption}
\usepackage{subcaption}
\captionsetup{font=small,labelfont=bf}
% Đặt ảnh theo cấu trúc: outputs/ và outputs/
% Thay đổi path bên dưới nếu bạn tổ chức thư mục khác trong project
\graphicspath{{./}{../outputs/}}

% ── Hyperlinks ────────────────────────────────────────────────────────────────
\usepackage[hidelinks,colorlinks=true,linkcolor=hust,urlcolor=hust]{hyperref}

% ── TOC depth ────────────────────────────────────────────────────────────────
\setcounter{tocdepth}{2}
\emergencystretch=3em
\sloppy

% ── Custom commands ──────────────────────────────────────────────────────────
\newcommand{\code}[1]{\begingroup\small\ttfamily\seqsplit{#1}\endgroup}
\newcommand{\tick}{$\checkmark$}
\newcommand{\cross}{$\times$}
\newcommand{\riskscore}{\text{RiskScore}}

% ─────────────────────────────────────────────────────────────────────────────
\begin{document}

% ══════════════════════════════════════════════════════════════════════════════
%  TRANG BÃŒA
% ══════════════════════════════════════════════════════════════════════════════
\newgeometry{top=1.6cm,bottom=1.6cm,left=2.2cm,right=2.2cm}
\begin{titlepage}
  \thispagestyle{empty}
  \centering
  \color{black}

  {\large\bfseries ĐẠI HỌC BÁCH KHOA HÀ NỘI\par}
  \vspace{0.15cm}
  {\large\bfseries TRƯỜNG CÔNG NGHỆ THÔNG TIN VÀ TRUYỀN THÔNG\par}

  \vspace{0.65cm}
  \rule{0.32\textwidth}{0.4pt}

  \vspace{0.9cm}
  {\LARGE\bfseries BÁO CÁO BÀI TẬP LỚN\par}
  \vspace{0.35cm}
  {\Large\bfseries NHẬP MÔN HỌC MÁY VÀ KHAI PHÁ DỮ LIỆU\par}

  \vspace{1.05cm}
  \includegraphics[width=0.48\textwidth]{logo-soict-hust-1.png}

  \vspace{0.95cm}
  {\large\bfseries TÊN ĐỀ TÀI\par}
  \vspace{0.35cm}
  \begin{minipage}{0.9\textwidth}
    \centering
    {\large\bfseries
    ỨNG DỤNG MACHINE LEARNING TRONG DỰ ĐOÁN VÀ PHÂN TẦNG RỦI RO KHÁCH HÀNG\\
    SO SÁNH HAI BÀI TOÁN: CUSTOMER CHURN VÀ PASSENGER SATISFACTION\par}
  \end{minipage}

  \vspace{0.95cm}
  \begin{fittabular}{@{}rl@{}}
    \textbf{Giảng viên hướng dẫn:} & Ngô Văn Linh \\[0.45em]
    \textbf{Sinh viên thực hiện:}  & Vũ Quang Nhật -- 20235800 \\
                                  & Phạm Tiến Đại -- 20225799 \\
                                  & Hoàng Nguyễn Phong -- 20215112 \\[0.45em]
    \textbf{Nhóm:}                 & 1
  \end{fittabular}

  \vfill
  {\itshape Hà Nội, tháng 6 năm 2026\par}
\end{titlepage}
\restoregeometry

% ══════════════════════════════════════════════════════════════════════════════
%  HEADER META

% ══════════════════════════════════════════════════════════════════════════════
%  HEADER META
% ══════════════════════════════════════════════════════════════════════════════
\begin{center}
  {\bfseries Nhóm thực hiện:} Vũ Quang Nhật · Phạm Tiến Đại · Hoàng Nguyễn Phong\\
  {\bfseries Phạm vi:} E-Commerce Customer Churn Prediction\\
  {\bfseries Công nghệ:} Python · XGBoost · KMeans · SHAP · Kaplan-Meier · Optuna
\end{center}

\rule{\textwidth}{0.4pt}

% ══════════════════════════════════════════════════════════════════════════════
%  MỤC LỤC
% ══════════════════════════════════════════════════════════════════════════════
\newpage
\tableofcontents
\newpage

% ══════════════════════════════════════════════════════════════════════════════
\section{Phát Biểu Bài Toán}
% ══════════════════════════════════════════════════════════════════════════════

\subsection{E-Commerce Customer Churn Prediction}

\textbf{Bối cảnh:}
Trong ngành thương mại điện tử, chi phí giữ chân một khách hàng cũ thấp hơn đáng kể so với thu
hút khách mới. Việc xác định sớm khách hàng có khả năng rời bỏ (\textit{churn}) giúp doanh nghiệp
can thiệp kịp thời và phân bổ nguồn lực CSKH hiệu quả hơn.

\textbf{Bài toán:}
Xây dựng pipeline dự đoán xác suất churn cho từng khách hàng, kết hợp phân nhóm hành vi
(clustering) và phân tích vòng đời (survival analysis), để xuất ra một \textbf{Risk Score cá nhân
hóa} giúp đội ngũ CSKH ưu tiên can thiệp đúng đối tượng, đúng thời điểm.

\textbf{Input:} Dữ liệu hành vi mua sắm của 5\,630 khách hàng (lịch sử đơn hàng, phàn nàn, thiết
bị, cashback, tenure).

\textbf{Output:} $\text{RiskScore} \in [0,1]$ $\times$ $\text{Cluster} \in \{0,1,2,3,4\}$ $\times$
Survival Context $\to$ Ma trận can thiệp ưu tiên (P1--P5).

\subsection*{Mục tiêu phân tích}

\begin{center}
\begin{tabularx}{\textwidth}{lX}
  \toprule
  \textbf{Thành phần} & \textbf{Mô tả} \\
  \midrule
  Loại bài toán      & Binary classification: dự đoán \code{Churn} (0/1) \\
  Đơn vị phân tích   & 1 dòng tương ứng 1 khách hàng trong dữ liệu thương mại điện tử \\
  Phân khúc hành vi  & RFM segmentation và KMeans để mô tả nhóm khách hàng \\
  Phân tích vòng đời & Kaplan-Meier và log-rank test dựa trên \code{Tenure} và \code{Churn} \\
  Đầu ra cuối        & Risk Score, cluster hành vi và ma trận ưu tiên can thiệp P1--P5 \\
  Vấn đề dữ liệu     & Mất cân bằng lớp churn (16.8\% churn, 83.2\% không churn) và một số cột có missing value \\
  \bottomrule
\end{tabularx}
\end{center}

% ══════════════════════════════════════════════════════════════════════════════
\section{Cơ Sở Lý Thuyết}
\label{sec:theory}
% ══════════════════════════════════════════════════════════════════════════════

Mục này trình bày nền tảng lý thuyết và công thức toán học của các kỹ thuật được sử dụng trong
pipeline E-Commerce Churn Alert. Nội dung được nhóm theo các giai đoạn chính của một dự án học máy:
\textit{K láng giềng/KNN}, \textit{phân tích thống kê (EDA)}, \textit{phân khúc khách hàng (RFM)},
\textit{phân cụm}, \textit{mô hình phân loại \& tinh chỉnh}, và \textit{diễn giải/hậu kiểm}.

\subsection{K Láng Giềng Gần Nhất (KNN)}

\subsubsection{Nguyên lý KNN}

KNN (\textit{K-Nearest Neighbors}) là một phương pháp học dựa trên mẫu
(\textit{instance-based learning}). Thay vì học một hàm tham số rõ ràng trong giai đoạn huấn luyện,
KNN lưu lại dữ liệu đã biết và khi cần dự đoán cho một điểm mới $x$, thuật toán tìm $k$ điểm gần
nhất với $x$ trong không gian đặc trưng. Nhãn hoặc giá trị của $x$ được suy ra từ các láng giềng
này.

Với bài toán phân loại, dự đoán thường được xác định bằng bỏ phiếu đa số:
\[
  \hat{y}(x) = \arg\max_{c \in \mathcal{C}}
  \sum_{i \in \mathcal{N}_k(x)} \mathbf{1}(y_i = c)
\]
trong đó $\mathcal{N}_k(x)$ là tập $k$ láng giềng gần nhất của $x$, $\mathcal{C}$ là tập nhãn và
$\mathbf{1}(\cdot)$ là hàm chỉ báo. Nếu dùng trọng số theo khoảng cách, láng giềng gần hơn có ảnh
hưởng lớn hơn:
\[
  \hat{y}(x) = \arg\max_{c \in \mathcal{C}}
  \sum_{i \in \mathcal{N}_k(x)} w_i\,\mathbf{1}(y_i = c),
  \qquad
  w_i = \frac{1}{d(x,x_i)+\epsilon}
\]

\subsubsection{Khoảng cách trong KNN}

KNN phụ thuộc mạnh vào cách đo khoảng cách. Với dữ liệu số đã được chuẩn hóa, khoảng cách Euclid là
lựa chọn phổ biến:
\[
  d(x_i, x_l) = \sqrt{\sum_{m=1}^{p}(x_{i,m} - x_{l,m})^2}
\]
Trong trường hợp có nhiều đặc trưng khác thang đo, các biến có giá trị lớn sẽ lấn át biến nhỏ hơn.
Vì vậy, với các mô hình hoặc bước xử lý dựa trên khoảng cách, dữ liệu cần được scale trước khi tính
khoảng cách. Trong pipeline e-commerce, điều này đặc biệt quan trọng với các biến như
\code{CashbackAmount}, \code{WarehouseToHome}, \code{OrderCount} và \code{Tenure}.

Giá trị $k$ điều khiển độ trơn của mô hình: $k$ nhỏ làm mô hình nhạy với nhiễu và outlier, còn $k$
lớn làm dự đoán ổn định hơn nhưng có thể làm mất ranh giới cục bộ. Do đó, $k$ thường được chọn bằng
thực nghiệm hoặc cross-validation. Trong bài toán này, $k=5$ được dùng cho bước điền khuyết vì cân
bằng giữa tính cục bộ và độ ổn định.

\subsubsection{KNN Imputer (Điền khuyết bằng K láng giềng gần nhất)}

Với mỗi mẫu có giá trị thiếu tại đặc trưng $j$, KNN Imputer tìm $k$ láng giềng gần nhất trong không
gian các đặc trưng \emph{không thiếu}, dựa trên khoảng cách Euclid có trọng số theo số chiều quan
sát được:
\[
  d(x_i, x_l) = \sqrt{\dfrac{\#\text{cols}}{\#\text{cols hiện diện ở cả hai mẫu}}
  \sum_{m \in \text{cols hiện diện}} (x_{i,m} - x_{l,m})^2}
\]
Giá trị điền khuyết là trung bình của $k$ láng giềng gần nhất:
\[
  \hat{x}_{i,j} = \dfrac{1}{k}\sum_{l \in \mathcal{N}_k(i)} x_{l,j}
\]
Nhóm dùng $k=5$ vì giữ được cấu trúc tương quan giữa các đặc trưng (so với mean/median imputation
chỉ điền một hằng số duy nhất, làm méo phương sai và tương quan).

Trong dữ liệu e-commerce, các cột như \code{Tenure}, \code{WarehouseToHome}, \code{HourSpendOnApp},
\code{OrderCount}, \code{CouponUsed}, \code{OrderAmountHikeFromlastYear} và
\code{DaySinceLastOrder} có tỉ lệ thiếu khoảng 4--5\%. Nếu thay toàn bộ bằng trung vị, mọi khách
hàng bị thiếu cùng một cột sẽ nhận cùng một giá trị, làm giảm khác biệt hành vi giữa các khách hàng.
KNN Imputer tận dụng các khách hàng có hồ sơ gần giống nhau để ước lượng giá trị thiếu, nhờ đó phù
hợp hơn với bài toán churn nơi các biến hành vi thường có tương quan với nhau.

\subsection{Các Kỹ Thuật Xử Lý Dữ Liệu Bổ Trợ}

\subsubsection{Phương pháp IQR --- Phát hiện \& Xử lý Outlier}

Khoảng liên phân vị (Interquartile Range):
\[
  IQR = Q_3 - Q_1
\]
trong đó $Q_1, Q_3$ là phân vị thứ 25\% và 75\%. Một điểm dữ liệu $x$ được xem là outlier nếu:
\[
  x < Q_1 - 1.5 \cdot IQR \quad \text{hoặc} \quad x > Q_3 + 1.5 \cdot IQR
\]
Thay vì loại bỏ (drop), nhóm áp dụng \textbf{clipping} (winsorizing): mọi giá trị vượt ngưỡng được
``ép'' về đúng biên:
\[
  x_{\text{clip}} = \min\big(\max(x,\, Q_1 - 1.5\,IQR),\; Q_3 + 1.5\,IQR\big)
\]
giúp giữ nguyên số lượng mẫu, tránh mất thông tin của các dòng khác trong cùng bản ghi.

\subsubsection{Min-Max Scaling}

Chuẩn hóa tuyến tính đưa mọi đặc trưng về cùng khoảng $[0,1]$, cần thiết cho các mô hình dựa trên
khoảng cách hoặc gradient (Logistic Regression):
\[
  x' = \dfrac{x - x_{\min}}{x_{\max} - x_{\min}}
\]
Khác với Z-score, MinMax không giả định phân phối Gauss và bảo toàn hình dạng phân phối gốc, phù
hợp sau khi đã clip outlier bằng IQR.

\subsubsection{One-Hot Encoding (OHE)}

Với biến hạng mục $X$ có $c$ giá trị duy nhất $\{v_1,\dots,v_c\}$, OHE tạo $c$ (hoặc $c-1$, tránh
multicollinearity) cột nhị phân:
\[
  X^{(v_k)}_i =
  \begin{cases}
    1 & \text{nếu } X_i = v_k \\
    0 & \text{ngược lại}
  \end{cases}
\]
Áp dụng cho các biến không có thứ tự tự nhiên trong dữ liệu e-commerce như
\code{PreferedOrderCat}, \code{PreferredPaymentMode}, \code{PreferredLoginDevice},
\code{MaritalStatus} và \code{Gender}. Với biến dạng nhóm nhưng có ý nghĩa mức độ như
\code{CityTier}, nhóm chuyển kiểu dữ liệu phù hợp trước khi đưa vào pipeline để tránh mô hình hiểu
sai quan hệ số học.

\subsubsection{SMOTE --- Synthetic Minority Over-sampling Technique}

Với mỗi mẫu $x_i$ thuộc lớp thiểu số, SMOTE chọn ngẫu nhiên một láng giềng gần nhất $x_{zi}$
(trong cùng lớp) và sinh mẫu tổng hợp bằng nội suy tuyến tính:
\[
  x_{\text{new}} = x_i + \lambda \cdot (x_{zi} - x_i), \qquad \lambda \sim \mathcal{U}(0,1)
\]
Việc này tạo thêm mẫu lớp churn (16.8\%) trên đường nối hai điểm thật trong không gian đặc trưng,
khác với random oversampling (chỉ lặp lại mẫu cũ, dễ overfit). SMOTE chỉ áp dụng \textbf{sau khi
chia train/test} và chỉ trên \textbf{tập train} để tránh leakage sang tập kiểm tra.

\subsubsection{Chi-bình phương ($\chi^2$) --- Lựa chọn đặc trưng}

Kiểm định mức độ phụ thuộc giữa đặc trưng (đã rời rạc hóa/không âm) và nhãn:
\[
  \chi^2 = \sum_{i=1}^{r}\sum_{j=1}^{c} \dfrac{(O_{ij} - E_{ij})^2}{E_{ij}}
\]
với $O_{ij}$ là tần suất quan sát, $E_{ij}$ là tần suất kỳ vọng nếu hai biến độc lập. Đặc trưng có
$p\text{-value} < 0.05$ (tương đương $\chi^2$ lớn) được giữ lại trong Linear Path --- đây là một
phép kiểm định thống kê, không phải khoảng cách, nên chỉ dùng được trên dữ liệu không âm (sau
MinMaxScale).

\subsection{Phân Tích Thống Kê Hai Biến (Bivariate EDA)}

\subsubsection{Hệ số tương quan Pearson}

Đo mức độ tương quan \emph{tuyến tính} giữa hai biến liên tục $X, Y$:
\[
  r = \dfrac{\sum_{i}(x_i - \bar{x})(y_i - \bar{y})}{\sqrt{\sum_i (x_i-\bar{x})^2}\sqrt{\sum_i (y_i-\bar{y})^2}}, \qquad r \in [-1, 1]
\]
Dùng để xây heatmap tương quan và phát hiện multicollinearity giữa các biến số trong dữ liệu
e-commerce. Ví dụ, nếu hai biến hành vi mua sắm có tương quan quá cao, nhóm sẽ cân nhắc giữ lại
biến có ý nghĩa nghiệp vụ rõ hơn để giảm nhiễu và giúp mô hình ổn định.

\subsubsection{Tương quan Point-Biserial}

Trường hợp đặc biệt của Pearson khi một biến là liên tục ($Y$) và biến còn lại là nhị phân ($X \in
\{0,1\}$, ví dụ \code{Churn}):
\[
  r_{pb} = \dfrac{\bar{Y}_1 - \bar{Y}_0}{s_Y}\sqrt{p\,(1-p)}
\]
trong đó $\bar{Y}_1,\bar{Y}_0$ là trung bình $Y$ tại nhóm $X=1$ và $X=0$, $s_Y$ là độ lệch chuẩn
của $Y$, $p$ là tỉ lệ $X=1$. Dùng để xếp hạng các đặc trưng số ảnh hưởng đến nhãn nhị phân
(\code{Tenure}: $r=-0.35$).

\subsubsection{Hệ số Cramér's V}

Đo mức độ liên kết giữa hai biến \emph{hạng mục}, dựa trên thống kê $\chi^2$ đã chuẩn hóa về
$[0,1]$:
\[
  V = \sqrt{\dfrac{\chi^2 / n}{\min(r-1,\, c-1)}}
\]
với $n$ là tổng số quan sát, $r, c$ là số hạng mục của hai biến. $V \to 0$: độc lập; $V \to 1$:
liên kết hoàn toàn. Ví dụ \code{Complain} có $V \approx 0.33$ với \code{Churn} --- liên kết mạnh
nhất trong các biến hạng mục.

\subsection{RFM Segmentation}

Mô hình RFM chấm điểm khách hàng theo 3 chiều hành vi giao dịch, sử dụng phân vị (\textbf{quintile,
NTILE(5)}) để chia mỗi chiều thành 5 nhóm điểm $1$--$5$:

\begin{itemize}
  \item \textbf{Recency (R):} $R_i = 5 - \text{quintile rank}(\code{DaySinceLastOrder}_i)$ --- mua
    gần đây hơn $\to$ điểm cao hơn.
  \item \textbf{Frequency (F):} $F_i = \text{quintile rank}(\code{OrderCount}_i)$ --- mua nhiều
    hơn $\to$ điểm cao hơn.
  \item \textbf{Monetary (M):} $M_i = \text{quintile rank}(\code{CashbackAmount}_i)$ --- giá trị
    cao hơn $\to$ điểm cao hơn.
\end{itemize}

Hàm \texttt{NTILE(5)} (SQL window function) chia một cột đã sắp xếp thành 5 phần có số lượng phần
tử gần bằng nhau:
\[
  \text{NTILE}(k)\text{ tại hàng } i = \left\lceil \dfrac{i \cdot k}{n} \right\rceil, \quad i = 1,\dots,n
\]
\code{rfm\_total} $= R + F + M \in [3,15]$ được dùng làm đặc trưng input cho mô hình ML, còn các
quy tắc ngưỡng (ví dụ $F\geq 4, M\geq 4 \to$ Champions) định nghĩa 5 \emph{segment} mô tả ở
Bảng~3.

\subsection{Phân Cụm (Clustering)}

\subsubsection{K-Means}

Thuật toán phân $n$ điểm dữ liệu thành $K$ cụm bằng cách tối thiểu hóa tổng bình phương khoảng cách
trong cụm (Within-Cluster Sum of Squares):
\[
  \arg\min_{S} \sum_{k=1}^{K} \sum_{x \in S_k} \lVert x - \mu_k \rVert^2
\]
với $\mu_k = \frac{1}{|S_k|}\sum_{x \in S_k} x$ là tâm cụm. Thuật toán Lloyd lặp 2 bước đến khi hội
tụ:
\begin{enumerate}
  \item \textbf{Assignment:} gán mỗi điểm vào cụm có tâm gần nhất: $\;c_i = \arg\min_k \lVert x_i - \mu_k\rVert^2$.
  \item \textbf{Update:} cập nhật lại tâm cụm là trung bình các điểm vừa gán.
\end{enumerate}
KMeans được \textbf{fit chỉ trên dữ liệu train (pre-SMOTE)} để tránh leakage; tập test dùng
\code{.predict()} với tâm cụm đã học. Hai dạng đặc trưng phái sinh từ KMeans:
\begin{itemize}
  \item \textbf{Distance features:} $d_k(x) = \lVert x - \mu_k \rVert$ cho mỗi $k=1,\dots,K$.
  \item \textbf{Cluster ID:} biến hạng mục $c(x) = \arg\min_k \lVert x - \mu_k\rVert$.
\end{itemize}

\subsubsection{Phương pháp Elbow}

Vẽ đường \textbf{WCSS} (hay \textit{inertia}) theo số cụm $K$:
\[
  \text{WCSS}(K) = \sum_{k=1}^{K}\sum_{x \in S_k} \lVert x - \mu_k \rVert^2
\]
$K$ tối ưu được chọn tại ``khuỷu tay'' (elbow) --- điểm mà việc tăng $K$ không còn giảm WCSS đáng kể.

\subsubsection{Hệ số Silhouette}

Đo chất lượng phân cụm cho từng điểm $i$:
\[
  s(i) = \dfrac{b(i) - a(i)}{\max\big(a(i), b(i)\big)}, \qquad s(i) \in [-1, 1]
\]
với $a(i)$ là khoảng cách trung bình từ $i$ đến các điểm cùng cụm, $b(i)$ là khoảng cách trung bình
nhỏ nhất từ $i$ đến các điểm của cụm khác gần nhất. Silhouette score toàn cục là trung bình
$s(i)$ trên toàn bộ mẫu. $s \to 1$: cụm tách biệt rõ; $s \to 0$: cụm chồng lấp; $s < 0$: gán sai
cụm. Với dữ liệu E-Commerce, silhouette $=0.1183$ tại $K=5$ cho thấy ranh giới cụm chưa thật sắc
nét, nhưng các cụm vẫn hữu ích để mô tả persona và bổ sung ngữ cảnh cho Risk Score.

\subsection{Mô Hình Phân Loại}

\subsubsection{Logistic Regression}

Mô hình tuyến tính cho xác suất nhãn nhị phân, dùng hàm sigmoid để nén $\mathbb{R} \to [0,1]$:
\[
  P(y=1\mid x) = \sigma(w^\top x + b) = \dfrac{1}{1+e^{-(w^\top x + b)}}
\]
Tham số $w, b$ được học bằng cực đại hóa log-likelihood, tương đương tối thiểu hóa \textbf{Binary
Cross-Entropy}:
\[
  \mathcal{L}(w,b) = -\dfrac{1}{n}\sum_{i=1}^n \Big[y_i \log \hat{y}_i + (1-y_i)\log(1-\hat{y}_i)\Big] + \lambda \lVert w \rVert_2^2
\]
với $\lambda \lVert w\rVert_2^2$ là $L_2$ regularization chống overfit. Đây là baseline tuyến
tính (LR-Base, LR-Dist) trong pipeline e-commerce.

\subsubsection{Decision Tree}

Cây quyết định chia đệ quy không gian đặc trưng bằng cách chọn ngưỡng $(j,\tau)$ tối thiểu hóa độ
``hỗn loạn'' (impurity) của hai node con. Với \textbf{Gini Impurity}:
\[
  G(t) = 1 - \sum_{k} p_k(t)^2
\]
hoặc \textbf{Entropy}:
\[
  H(t) = -\sum_k p_k(t)\log_2 p_k(t)
\]
với $p_k(t)$ là tỉ lệ lớp $k$ tại node $t$. Tại mỗi node, thuật toán chọn split tối đa hóa
\textbf{Information Gain}:
\[
  IG(t, j, \tau) = G(t) - \dfrac{n_L}{n_t}G(t_L) - \dfrac{n_R}{n_t}G(t_R)
\]
DT-Clust là baseline phi tuyến đơn giản, dễ diễn giải nhưng dễ overfit (thể hiện qua Accuracy/F1
thấp hơn rõ rệt so với ensemble trong Bảng kết quả).

\subsubsection{Random Forest}

Kỹ thuật \textbf{Bagging} (Bootstrap Aggregating): huấn luyện $B$ cây quyết định độc lập trên các
mẫu bootstrap (lấy lại mẫu có hoàn lại) của tập train, mỗi cây chỉ xét một tập con ngẫu nhiên đặc
trưng tại mỗi split (giảm tương quan giữa các cây). Dự đoán cuối là trung bình (regression) hoặc
bầu chọn đa số/trung bình xác suất (classification):
\[
  \hat{P}(y=1\mid x) = \dfrac{1}{B}\sum_{b=1}^{B} P_b(y=1\mid x)
\]
Giảm phương sai (variance) so với một cây đơn lẻ nhờ lấy trung bình nhiều mô hình yếu tương quan
thấp $\to$ RF-Clust đạt ROC-AUC rất cao (0.9975 ở E-Commerce) nhưng threshold cao bất thường (0.390)
cho thấy hiện tượng \textit{under-calibration}.

\subsubsection{XGBoost --- Extreme Gradient Boosting}

Mô hình \textbf{Gradient Boosting} xây dựng tuần tự $T$ cây, mỗi cây sau học để sửa lỗi (residual
theo gradient) của tổng các cây trước:
\[
  \hat{y}_i^{(t)} = \hat{y}_i^{(t-1)} + \eta \cdot f_t(x_i)
\]
với $\eta$ là learning rate. Hàm mục tiêu tại bước $t$ gồm hàm mất mát và regularization trên độ
phức tạp cây:
\[
  \mathcal{L}^{(t)} = \sum_{i=1}^n \ell\big(y_i, \hat{y}_i^{(t-1)} + f_t(x_i)\big) + \Omega(f_t),
  \qquad \Omega(f) = \gamma T + \tfrac{1}{2}\lambda \sum_{j=1}^{T} w_j^2
\]
với $T$ là số lá, $w_j$ trọng số lá $j$, $\gamma,\lambda$ là hệ số regularization. XGBoost xấp xỉ
$\mathcal{L}^{(t)}$ bằng khai triển Taylor bậc 2 (dùng gradient $g_i$ và Hessian $h_i$ của hàm mất
mát), cho phép tính \textbf{Gain} khi xét một split ứng viên:
\[
  \text{Gain} = \tfrac{1}{2}\left[\dfrac{(\sum_{i\in L} g_i)^2}{\sum_{i\in L} h_i + \lambda} +
  \dfrac{(\sum_{i\in R} g_i)^2}{\sum_{i\in R} h_i + \lambda} -
  \dfrac{(\sum_{i\in L\cup R} g_i)^2}{\sum_{i\in L\cup R} h_i + \lambda}\right] - \gamma
\]
XGBoost xử lý quan hệ phi tuyến và tương tác đặc trưng rất tốt, phù hợp với dữ liệu churn nơi hành
vi của khách hàng thường không tác động độc lập. Trong pipeline e-commerce, XGB-Clust được chọn làm
model production nhờ cân bằng tốt giữa ROC-AUC, F1, khả năng calibrate và khả năng tận dụng thêm
ngữ cảnh phân cụm.

\subsection{Tinh Chỉnh Siêu Tham Số --- Optuna}

Optuna dùng thuật toán Bayesian Optimization theo cơ chế \textbf{TPE (Tree-structured Parzen
Estimator)}: xây hai phân phối mật độ xác suất $\ell(x)$ (cho các thử nghiệm ``tốt'') và $g(x)$ (cho
các thử nghiệm ``kém'') trên không gian siêu tham số, rồi chọn điểm thử tiếp theo tối đa hóa tỉ số:
\[
  x_{\text{next}} = \arg\max_x \dfrac{\ell(x)}{g(x)}
\]
Khác với Grid/Random Search (thử ngẫu nhiên/độc lập), TPE học dần từ lịch sử các trial trước để tập
trung tìm kiếm vào vùng siêu tham số nhiều khả năng tối ưu, hiệu quả hơn khi số trial hạn chế. Trong
pipeline e-commerce, nhóm dùng 40 trials để tinh chỉnh mô hình mà vẫn kiểm soát thời gian chạy.

\subsection{Hiệu Chỉnh Xác Suất --- CalibratedClassifierCV}

Một mô hình ``tốt'' về phân loại (AUC cao) không đồng nghĩa xác suất đầu ra $\hat{p}$ phản ánh đúng
tần suất thật. \textbf{Platt Scaling} học một phép biến đổi sigmoid trên điểm số thô $f(x)$ của mô
hình:
\[
  P_{\text{calibrated}}(y=1\mid x) = \dfrac{1}{1+\exp(A\cdot f(x) + B)}
\]
với $A,B$ học bằng cross-validation trên một phần dữ liệu giữ lại (\code{cv=5}). Mô hình được
calibrate tốt khi đường \textbf{Calibration Curve} (xác suất dự đoán vs tần suất thực tế quan
sát) gần với đường chéo $y=x$ --- đây là tiêu chí nhóm dùng để giải thích vì sao chọn XGB-Clust dù
RF-Clust có AUC cao hơn.

\subsection{Chỉ Số Đánh Giá Mô Hình}

Từ ma trận nhầm lẫn (Confusion Matrix) với $TP, FP, FN, TN$:
\[
  \text{Accuracy} = \dfrac{TP+TN}{TP+FP+FN+TN}, \qquad
  \text{Precision} = \dfrac{TP}{TP+FP}, \qquad
  \text{Recall} = \dfrac{TP}{TP+FN}
\]
\[
  F_1 = 2\cdot\dfrac{\text{Precision}\cdot\text{Recall}}{\text{Precision}+\text{Recall}}
\]
\textbf{ROC-AUC:} diện tích dưới đường cong ROC (\textit{True Positive Rate} vs \textit{False
Positive Rate} khi quét ngưỡng phân loại từ $0$ đến $1$):
\[
  TPR = \dfrac{TP}{TP+FN}, \qquad FPR = \dfrac{FP}{FP+TN}, \qquad \text{AUC} = \int_0^1 TPR\, d(FPR)
\]
\textbf{PR-AUC:} diện tích dưới đường cong Precision--Recall, nhạy hơn ROC-AUC khi lớp dương là
thiểu số (phù hợp với churn 16.8\%):
\[
  \text{PR-AUC} = \int_0^1 \text{Precision}(r)\, dr
\]

\subsection{Diễn Giải Mô Hình --- SHAP (SHapley Additive exPlanations)}

SHAP áp dụng \textbf{giá trị Shapley} từ lý thuyết trò chơi hợp tác để phân bổ ``công bằng'' đóng góp
của mỗi đặc trưng $j$ vào một dự đoán cụ thể $f(x)$. Với $N$ là tập toàn bộ đặc trưng:
\[
  \phi_j = \sum_{S \subseteq N \setminus \{j\}} \dfrac{|S|!\,(|N|-|S|-1)!}{|N|!}
  \Big[f(S \cup \{j\}) - f(S)\Big]
\]
trong đó $f(S)$ là giá trị dự đoán kỳ vọng khi chỉ biết các đặc trưng trong tập con $S$. Tính chất
quan trọng nhất là \textbf{tính cộng (additivity)}: tổng các $\phi_j$ cộng với giá trị nền (baseline)
$\phi_0 = E[f(x)]$ đúng bằng dự đoán thật:
\[
  f(x) = \phi_0 + \sum_{j=1}^{|N|} \phi_j
\]
Vì tính $2^{|N|}$ tập con là không khả thi, nhóm dùng \textbf{TreeSHAP} --- thuật toán xấp xỉ chính
xác (exact) giá trị Shapley cho mô hình dạng cây (XGBoost) trong thời gian đa thức bằng cách truy
vết đường đi qua cây. Với bài toán churn, SHAP giúp xác định các đặc trưng đẩy xác suất rời bỏ lên
cao như \code{Tenure}, \code{Complain}, \code{DaySinceLastOrder} hoặc các biến liên quan đến giá trị
mua hàng.

\subsection{Phân Tích Sống Còn --- Kaplan-Meier \& Log-rank Test}

\subsubsection{Kaplan-Meier Estimator}

Ước lượng \textbf{hàm sống còn} $S(t)$ --- xác suất một khách hàng ``sống sót'' (chưa churn) qua thời
điểm $t$ --- theo phương pháp non-parametric, xử lý đúng dữ liệu \textit{censored} (khách chưa
churn tại thời điểm quan sát cuối):
\[
  \hat{S}(t) = \prod_{t_i \leq t} \left(1 - \dfrac{d_i}{n_i}\right)
\]
với $t_i$ là các thời điểm xảy ra event (churn), $d_i$ là số event tại $t_i$, $n_i$ là số đối
tượng ``có nguy cơ'' (at risk) ngay trước $t_i$. Trong pipeline, \code{duration = Tenure},
\code{event = Churn}.

\subsubsection{Log-rank Test}

Kiểm định giả thuyết $H_0$: các nhóm (cluster) có hàm sống còn giống nhau, dựa trên so sánh số
event quan sát $O_{k}$ và kỳ vọng $E_k$ tại mỗi nhóm $k$ qua tất cả các thời điểm event:
\[
  \chi^2 = \sum_{k} \dfrac{(O_k - E_k)^2}{E_k} \;\sim\; \chi^2_{(\text{df}=K-1)}
\]
$p$-value nhỏ ($8.69\times 10^{-9}$ trong báo cáo) bác bỏ $H_0$ --- xác nhận 4 cluster có
\textbf{tốc độ rời bỏ} khác biệt có ý nghĩa thống kê, bổ sung chiều thông tin ``khi nào'' bên cạnh
``có churn hay không'' của mô hình phân loại.

% ══════════════════════════════════════════════════════════════════════════════
\section{Data Dictionary}
% ══════════════════════════════════════════════════════════════════════════════

\subsection{E-Commerce Customer Churn Dataset}

\textbf{Nguồn:} Kaggle --- E-Commerce Customer Churn Dataset\\
\textbf{Kích thước:} 5\,630 dòng $\times$ 20 cột · Churn rate: \textbf{16.8\%}

\begin{longtable}{llp{5cm}p{4.5cm}}
  \toprule
  \textbf{Cột} & \textbf{Kiểu} & \textbf{Mô tả} & \textbf{Ghi chú} \\
  \midrule
  \endhead
  \code{CustomerID}                 & int        & Định danh khách hàng & Unique key \\
  \code{Churn}                      & int (0/1)  & Đã churn trong kỳ? & \textbf{Target} \\
  \code{Tenure}                     & float      & Số tháng là thành viên & 4.69\% NaN $\to$ MAR/MNAR \\
  \code{CityTier}                   & cat        & Tier thành phố (1/2/3) & Chuyển về object \\
  \code{WarehouseToHome}            & float      & Khoảng cách kho $\to$ nhà (km) & 4.46\% NaN \\
  \code{HourSpendOnApp}             & float      & Giờ dùng app/tuần & 4.53\% NaN \\
  \code{NumberOfDeviceRegistered}   & int        & Số thiết bị đã đăng ký & \\
  \code{PreferedOrderCat}           & cat        & Danh mục mua sắm ưa thích & Mobile Phone / Laptop / Fashion / \ldots \\
  \code{SatisfactionScore}          & int        & Điểm hài lòng (1--5) & \\
  \code{MaritalStatus}              & cat        & Tình trạng hôn nhân & Single / Married / Divorced \\
  \code{NumberOfAddress}            & int        & Số địa chỉ giao hàng đã lưu & \\
  \code{Complain}                   & int (0/1)  & Có phàn nàn trong kỳ? & Liên kết mạnh nhất với churn \\
  \code{OrderAmountHikeFromlastYear}& float      & \% tăng giá trị đơn so với năm ngoái & 4.71\% NaN \\
  \code{CouponUsed}                 & float      & Số coupon đã dùng & 4.55\% NaN \\
  \code{OrderCount}                 & float      & Số đơn hàng trong kỳ & 4.58\% NaN \\
  \code{DaySinceLastOrder}          & float      & Số ngày kể từ đơn cuối & 5.45\% NaN $\to$ MCAR \\
  \code{CashbackAmount}             & float      & Tổng cashback nhận được (\$) & Feature quan trọng trong Risk Score \\
  \code{PreferredLoginDevice}       & cat        & Thiết bị đăng nhập ưa thích & Mobile Phone / Computer \\
  \code{PreferredPaymentMode}       & cat        & Phương thức thanh toán ưa thích & Credit Card / COD / UPI / \ldots \\
  \code{Gender}                     & cat        & Giới tính & Male / Female \\
  \bottomrule
\end{longtable}

\textbf{Composite features được tạo thêm:}

\begin{center}
\begin{tabularx}{\textwidth}{lXX}
  \toprule
  \textbf{Feature} & \textbf{Công thức} & \textbf{Ý nghĩa} \\
  \midrule
  \code{device\_per\_tenure}  & \code{NumberOfDeviceRegistered / (Tenure + 1)} & Mật độ thiết bị theo thời gian \\
  \code{cashback\_per\_order} & \code{CashbackAmount / (OrderCount + 1)} & Giá trị cashback mỗi đơn \\
  \code{inactivity\_ratio}    & \code{DaySinceLastOrder / (Tenure + 1)} & Tỷ lệ không hoạt động \\
  \bottomrule
\end{tabularx}
\end{center}

% ══════════════════════════════════════════════════════════════════════════════
\section{Phương Pháp}
% ══════════════════════════════════════════════════════════════════════════════

\subsection{Pipeline E-Commerce Churn Alert (9 Steps)}

\begin{lstlisting}
Raw Data (Excel 5,630 rows x 20 cols)
         |
         v
[Step 1]  Data Cleaning -- typo, dtype correction
         |
         v
[Step 2]  Missing Data Analysis -- MCAR/MAR/MNAR classification
         |
         v
[Step 3]  EDA -- Cramer's V, Point-Biserial, Pearson heatmap
         |
         v
[Step 4]  RFM Segmentation -- SQLite NTILE(5) window function
         |
    +----+-----------------------------+
    v                                  v
[Step 5] Tree Path              [Step 5b] Linear Path
  KNN Impute -> OHE -> SMOTE      IQR Clip -> OHE -> MinMaxScale
  -> KMeans FE (K=5)              -> KNN Impute -> Chi2 Selection
    |                                          |
    +------------------+-----------------------+
                       v
           [Step 6]  Modelling (6 models)
                       |
                       v
           [Step 7]  Cluster Profiling (5 Personas)
                       |
                       v
           [Step 8]  Risk Scoring
           RiskScore = 0.7 x P(churn) + 0.3 x cashback_pct
                       |
                       v
           [Step 9]  Kaplan-Meier Survival Analysis
                       |
                       v
    RiskScore + Cluster + Survival -> Action Matrix (P1-P5)
\end{lstlisting}

\subsubsection{Step 1 --- Data Cleaning}

Load Excel (sheet \code{E Comm}) $\to$ 5\,630 rows $\times$ 20 cột. Chuẩn hóa typo categorical:

\begin{center}
\begin{tabular}{ll}
  \toprule
  \textbf{Trước} & \textbf{Sau} \\
  \midrule
  \code{"Mobile Phone"}, \code{"Phone"} & \code{"Mobile"} \\
  \code{"CC"} & \code{"Credit Card"} \\
  \code{"COD"} & \code{"Cash on Delivery"} \\
  \bottomrule
\end{tabular}
\end{center}

\subsubsection{Step 2 --- Missing Data Analysis}

Phân loại MCAR/MAR/MNAR bằng cách so sánh churn rate giữa nhóm có và không có missing
($|\Delta\text{churn}| > 5\% \to$ MAR/MNAR):

\begin{center}
\begin{tabularx}{\textwidth}{lllX}
  \toprule
  \textbf{Cột} & \textbf{Missing} & \textbf{Loại} & \textbf{Xử lý} \\
  \midrule
  \code{DaySinceLastOrder}           & 5.45\% & MCAR     & KNN Imputer $k=5$ \\
  \code{OrderAmountHikeFromlastYear} & 4.71\% & MAR/MNAR & KNN Imputer $k=5$ \\
  \code{Tenure}                      & 4.69\% & MAR/MNAR & KNN Imputer $k=5$ \\
  \code{OrderCount}                  & 4.58\% & MAR/MNAR & KNN Imputer $k=5$ \\
  \code{CouponUsed}                  & 4.55\% & MAR/MNAR & KNN Imputer $k=5$ \\
  \code{HourSpendOnApp}              & 4.53\% & MAR/MNAR & KNN Imputer $k=5$ \\
  \code{WarehouseToHome}             & 4.46\% & MAR/MNAR & KNN Imputer $k=5$ \\
  \bottomrule
\end{tabularx}
\end{center}

\begin{quote}
  Phần lớn missing là MAR/MNAR $\to$ dùng \textbf{KNN Imputer ($k=5$)} thay vì mean/median để
  bảo toàn cấu trúc tương quan.
\end{quote}

\begin{figure}[htbp]
  \centering
  \includegraphics[width=0.85\textwidth]{02_missing_matrix.png}
  \caption{Ma trận missing — 7 cột thiếu dữ liệu (4.46--5.45\%), không xếp theo dòng/cột cố định
    $\to$ không phải lỗi thu thập hệ thống.}
  \label{fig:ecom-missing}
\end{figure}

\subsubsection{Step 3 --- EDA}

\begin{center}
\begin{tabularx}{\textwidth}{lX}
  \toprule
  \textbf{Phân tích} & \textbf{Phát hiện chính} \\
  \midrule
  Phân bố nhãn          & Mất cân bằng nặng: 83.2\% retained vs \textbf{16.8\% churn} $\to$ cần SMOTE \\
  Churn theo Tenure      & Churn cao nhất ở 0--3 tháng đầu ($\approx 35\%$), giảm dần theo thời gian \\
  Churn theo SatisfactionScore & Score 1 $\to$ churn $\approx 25\%$; Score 5 $\to$ churn $\approx 10\%$ \\
  Cramér's V             & \code{Complain} liên kết mạnh nhất với churn ($V \approx 0.33$) \\
  Point-Biserial         & Top features: \code{Tenure} ($r=-0.35$), \code{CashbackAmount}, \code{SatisfactionScore} \\
  \bottomrule
\end{tabularx}
\end{center}

\begin{figure}[htbp]
  \centering
  \begin{subfigure}[b]{0.48\textwidth}
    \includegraphics[width=\textwidth]{03_class_distribution.png}
    \caption{Phân bố nhãn Churn}
  \end{subfigure}
  \hfill
  \begin{subfigure}[b]{0.48\textwidth}
    \includegraphics[width=\textwidth]{03_pearson_heatmap.png}
    \caption{Pearson heatmap}
  \end{subfigure}

  \vspace{0.5em}

  \begin{subfigure}[b]{0.48\textwidth}
    \includegraphics[width=\textwidth]{03_cramers_v.png}
    \caption{Cramér's V — biến categorical}
  \end{subfigure}
  \hfill
  \begin{subfigure}[b]{0.48\textwidth}
    \includegraphics[width=\textwidth]{03_point_biserial.png}
    \caption{Point-Biserial correlation}
  \end{subfigure}
  \caption{EDA — E-Commerce Churn Dataset}
  \label{fig:ecom-eda}
\end{figure}

\begin{figure}[htbp]
  \centering
  \begin{subfigure}[b]{0.48\textwidth}
    \includegraphics[width=\textwidth]{03_tenure_satisfaction.png}
    \caption{Churn theo Tenure và SatisfactionScore}
  \end{subfigure}
  \hfill
  \begin{subfigure}[b]{0.48\textwidth}
    \includegraphics[width=\textwidth]{03_cashback_churn.png}
    \caption{Cashback Amount theo nhóm Churn}
  \end{subfigure}
  \caption{Khách churn nhận cashback trung bình thấp hơn rõ rệt — xác nhận vai trò của
    \code{CashbackAmount} trong Risk Score.}
  \label{fig:ecom-tenure-cashback}
\end{figure}

\subsubsection{Step 4 --- RFM Segmentation}

Dùng \textbf{SQLite in-memory + NTILE(5)} window function để chấm điểm R/F/M từ 1--5:

\begin{center}
\begin{tabular}{lll}
  \toprule
  \textbf{Chiều} & \textbf{Mapping} & \textbf{Hướng điểm} \\
  \midrule
  \textbf{R} (Recency)   & \code{DaySinceLastOrder} & Ít ngày $\to$ điểm cao \\
  \textbf{F} (Frequency) & \code{OrderCount}         & Nhiều đơn $\to$ điểm cao \\
  \textbf{M} (Monetary)  & \code{CashbackAmount}     & Nhiều cashback $\to$ điểm cao \\
  \bottomrule
\end{tabular}
\end{center}

\begin{center}
\begin{tabular}{lll}
  \toprule
  \textbf{Segment} & \textbf{Điều kiện} & \textbf{Số KH} \\
  \midrule
  Champions        & $F \geq 4$ và $M \geq 4$ & 1\,505 \\
  Recent Customers & $R \geq 4$ và $F < 3$   & 1\,319 \\
  Hibernating      & Còn lại                  & 1\,270 \\
  Loyal            & $F \geq 3$ và $M \geq 3$ & 1\,244 \\
  At Risk          & $R \leq 2$ và $F \geq 3$ & 292    \\
  \bottomrule
\end{tabular}
\end{center}

\begin{quote}
  RFM standalone ROC-AUC = \textbf{0.5503} --- yếu khi dùng độc lập, nhưng \code{rfm\_segment}
  và \code{rfm\_total} là input features quan trọng cho model chính.
\end{quote}

\begin{figure}[htbp]
  \centering
  \begin{subfigure}[b]{0.48\textwidth}
    \includegraphics[width=\textwidth]{04_rfm_churn_rate.png}
    \caption{Churn rate theo RFM Segment}
  \end{subfigure}
  \hfill
  \begin{subfigure}[b]{0.48\textwidth}
    \includegraphics[width=\textwidth]{04_rfm_segments_pie.png}
    \caption{Phân bố RFM Segment}
  \end{subfigure}
  \caption{Kết quả phân đoạn RFM — 5 segments từ SQLite NTILE(5)}
  \label{fig:ecom-rfm}
\end{figure}

\subsubsection{Step 5 --- Dual Preprocessing Pipeline}

\textbf{Tree Path:} KNN Impute $\to$ OHE $\to$ SMOTE (3\,746$\to$3\,746 each class) $\to$ KMeans FE
($K=5$, silhouette $= 0.1183$)
\begin{itemize}
  \item Strategy A: 4 khoảng cách centroid (\code{X\_train\_dist})
  \item Strategy B: 1 cột \code{kmeans\_cluster\_id} (\code{X\_train\_clust})
\end{itemize}

\textbf{Linear Path:} IQR clip ($\times 1.5$) $\to$ OHE $\to$ MinMaxScale $\to$ KNN Impute $\to$
Chi$^2$ ($p < 0.05$) $\to$ SMOTE
\begin{itemize}
  \item Base: \textbf{23/35 features} giữ lại
  \item Dist: \textbf{25/39 features} (thêm 4 KMeans distance features)
\end{itemize}

\begin{figure}[htbp]
  \centering
  \begin{subfigure}[b]{0.48\textwidth}
    \includegraphics[width=\textwidth]{05_elbow_silhouette.png}
    \caption{Elbow \& Silhouette — KMeans $K=5$}
  \end{subfigure}
  \hfill
  \begin{subfigure}[b]{0.48\textwidth}
    \includegraphics[width=\textwidth]{05b_chi2_scores.png}
    \caption{Chi$^2$ feature scores (Linear path)}
  \end{subfigure}
  \caption{Dual preprocessing — lựa chọn $K$ và feature selection}
  \label{fig:ecom-preproc}
\end{figure}

\subsubsection{Step 6 --- Modelling}

6 models song song: XGB-Base, XGB-Clust, RF-Clust, DT-Clust, LR-Base, LR-Dist.\\
XGB-Base tuned bằng \textbf{Optuna (40 trials, CV=5)}: \code{n\_estimators=463, max\_depth=8,
lr=0.0834, subsample=0.679}.\\
Tất cả được calibrate bằng \code{CalibratedClassifierCV(cv=5)}.

\begin{figure}[htbp]
  \centering
  \includegraphics[width=0.85\textwidth]{06b_model_comparison.png}
  \caption{So sánh 6 models — E-Commerce pipeline}
  \label{fig:ecom-model-cmp}
\end{figure}

\begin{figure}[htbp]
  \centering
  \begin{subfigure}[b]{0.48\textwidth}
    \includegraphics[width=\textwidth]{06_roc_curve.png}
    \caption{ROC Curve}
  \end{subfigure}
  \hfill
  \begin{subfigure}[b]{0.48\textwidth}
    \includegraphics[width=\textwidth]{06_confusion_matrix.png}
    \caption{Confusion Matrix}
  \end{subfigure}

  \vspace{0.5em}

  \begin{subfigure}[b]{0.48\textwidth}
    \includegraphics[width=\textwidth]{06_calibration_curve.png}
    \caption{Calibration Curve — xác nhận model calibrate tốt}
  \end{subfigure}
  \hfill
  \begin{subfigure}[b]{0.48\textwidth}
    \includegraphics[width=\textwidth]{06_shap_bar.png}
    \caption{SHAP Bar — Global Importance}
  \end{subfigure}
  \caption{Đánh giá model — E-Commerce Churn Alert}
  \label{fig:ecom-model-eval}
\end{figure}

\begin{figure}[htbp]
  \centering
  \includegraphics[width=0.85\textwidth]{06_shap_summary.png}
  \caption{SHAP Summary Plot — phân phối giá trị SHAP theo feature}
  \label{fig:ecom-shap-summary}
\end{figure}

\subsubsection{Step 7 --- Cluster Profiling}

KMeans $K=5$ fit trên 4\,504 pre-SMOTE training samples. Profile 5 personas từ test set (1\,126 KH)
--- xem Mục~\ref{sec:cluster-ecom}.

\subsubsection{Step 8 --- Risk Scoring}

\begin{lstlisting}
RiskScore = normalize( 0.7 x P(churn) + 0.3 x (CashbackAmount / OrderCount) )
\end{lstlisting}

Phân ngưỡng quintile: Very Low / Low / Medium / High / Critical (mỗi tier $\approx 20\%$).

\subsubsection{Step 9 --- Kaplan-Meier Survival Analysis}

\code{lifelines.KaplanMeierFitter} với \code{event=Churn}, \code{duration=Tenure}. Log-rank test
kiểm định sự khác biệt giữa 4 clusters --- xem Mục~\ref{sec:survival}.

% ══════════════════════════════════════════════════════════════════════════════
\section{Kết Quả}
% ══════════════════════════════════════════════════════════════════════════════

\subsection{Kết Quả E-Commerce Churn Alert}

\subsubsection{Model Comparison}

\begin{center}
\begin{tabular}{llllll}
  \toprule
  \textbf{Model} & \textbf{Accuracy} & \textbf{F1} & \textbf{ROC-AUC} & \textbf{PR-AUC} & \textbf{Threshold} \\
  \midrule
  \textbf{XGB-Base}     & 0.9591          & 0.8889          & 0.9799          & 0.9490          & 0.053 \\
  RF-Clust              & 0.9796          & 0.9421          & \textbf{0.9975} & \textbf{0.9869} & 0.390 \\
  \textbf{XGB-Clust} $\star$ & 0.9316 & 0.8230     & 0.9743          & 0.9340          & 0.075 \\
  DT-Clust              & 0.8055          & 0.6178          & 0.9197          & 0.7761          & 0.091 \\
  LR-Dist               & 0.8446          & 0.6479          & 0.9042          & 0.7374          & 0.501 \\
  LR-Base               & 0.8579          & 0.6507          & 0.8971          & 0.7311          & 0.552 \\
  \bottomrule
\end{tabular}
\end{center}

\textbf{$\star$ Production model: XGB-Clust} --- tích hợp cluster context, cân bằng tốt giữa
Precision (73.1\%) và Recall (94.2\%), ROC-AUC $= 0.974$.

\begin{quote}
  \textit{Lý do không chọn RF-Clust dù ROC-AUC cao nhất (0.9975): RF-Clust có threshold 0.390
  (cao hơn nhiều so với 0.053--0.075 của XGB), cho thấy mô hình kém calibrate và dễ miss nhóm
  churn probability thấp nhưng thực sự churn.}
\end{quote}

\subsubsection{Cluster Profiling --- 5 Personas}
\label{sec:cluster-ecom}

\begin{center}
\begin{tabular}{llllll}
  \toprule
  \textbf{Cluster} & \textbf{Persona} & \textbf{Tenure} & \textbf{Cashback} & \textbf{Churn Rate} & \textbf{Size} \\
  \midrule
  C0 & Khách Hàng Thân Thiết & 11.24 tháng & \$187.87 & 18\% & 134 KH \\
  C1 & Mua Sắm Thường Xuyên   & 9.91 tháng  & \$173.28 & 23\% & 224 KH \\
  C2 & Người Mua Thụ Động     & 8.52 tháng  & \$175.06 & 24\% & 261 KH \\
  C3 & Khách Hàng Ổn Định     & 10.73 tháng & \$177.55 & 12\% & 261 KH \\
  C4 & VIP Loyal Customer     & 11.12 tháng & \$173.98 & 9\%  & 246 KH \\
  \bottomrule
\end{tabular}
\end{center}

\begin{quote}
  \textbf{Quan sát C2:} Nhóm "Người Mua Thụ Động" (Cluster 2) có churn rate cao nhất (24\%) mặc dù
  Tenure không phải là thấp nhất. Đây là khách hàng ít tương tác, cần can thiệp bằng các offer cá nhân hóa.
  Ngược lại, C4 (VIP Loyal) có churn rate thấp nhất (9\%) --- nhóm này cần giữ chân bằng các quyền lợi exclusive.
\end{quote}

\begin{figure}[htbp]
  \centering
  \includegraphics[width=0.85\textwidth]{07_cluster_profiling.png}
  \caption{Cluster Profiling — 5 Personas E-Commerce}
  \label{fig:ecom-cluster}
\end{figure}

\subsubsection{Survival Analysis}
\label{sec:survival}

\textbf{Log-rank test:} $\chi^2 \approx 40.42$, \textbf{$p$-value $\approx 8.69 \times 10^{-9}$} $\to$ Rất có ý
nghĩa thống kê.

Cluster 2 có survival curve giảm nhanh nhất trong giai đoạn 0--6 tháng --- rời bỏ sớm trong vòng đời. Cluster 4 (VIP) thể hiện survival cao nhất.

\begin{figure}[htbp]
  \centering
  \includegraphics[width=0.85\textwidth]{08_km_by_cluster.png}
  \caption{Kaplan-Meier Survival Curve theo Cluster — Log-rank $p = 8.69 \times 10^{-9}$}
  \label{fig:ecom-km}
\end{figure}

\subsubsection{Risk Distribution \& Action Matrix}

\begin{center}
\begin{tabularx}{\textwidth}{lllX}
  \toprule
  \textbf{Risk Tier} & \textbf{Ngưỡng} & \textbf{Số KH} & \textbf{Hành động} \\
  \midrule
  Critical  & $> 0.4864$          & 225 & \textbf{P1 --- Can thiệp ngay lập tức} \\
  High      & $0.3438$--$0.4864$  & 225 & P2 --- Trong 24--48h \\
  Medium    & $0.2546$--$0.3438$  & 225 & P3 --- Tuần tới \\
  Low       & $0.1403$--$0.2546$  & 225 & P4 --- Tháng tới \\
  Very Low  & $\leq 0.1403$       & 226 & P5 --- Chăm sóc tiêu chuẩn \\
  \bottomrule
\end{tabularx}
\end{center}

% ══════════════════════════════════════════════════════════════════════════════
\section{Kết Luận}
% ══════════════════════════════════════════════════════════════════════════════

Pipeline E-Commerce Churn Alert xây dựng thành công hệ thống phân tầng rủi ro đa chiều cho 5\,630
khách hàng:

\begin{itemize}
  \item \textbf{Model XGB-Clust} đạt ROC-AUC $= 0.974$, F1 $= 0.823$ --- hiệu quả cao trong điều
    kiện dữ liệu mất cân bằng nặng (16.8\% churn).
  \item \textbf{5 clusters} phân tách 5 hành vi mua sắm khác nhau: từ VIP Loyal (C4: 9\% churn) đến
    Passive Buyers (C2: 24\% churn). Phân đoạn rõ ràng hỗ trợ CSKH trong targeting.
  \item \textbf{Kaplan-Meier} xác nhận sự khác biệt survival curve có ý nghĩa thống kê cao
    ($p \approx 8.69 \times 10^{-9}$), cung cấp thêm chiều ``tốc độ rời bỏ'' ngoài xác suất churn.
  \item \textbf{Action Matrix P1--P5} tích hợp cả 3 trụ cột (RiskScore + Cluster + Survival), cho
    phép CSKH ưu tiên chính xác và cá nhân hóa hành động theo từng nhóm hành vi.
\end{itemize}

\textbf{Nhóm nguy hiểm tuyệt đối:} \code{Critical Risk + Cluster 2 (Passive Buyer) + Tenure < 3 tháng} --- tất cả
dấu hiệu xấu đồng thời, cần can thiệp trong vài giờ.

% ══════════════════════════════════════════════════════════════════════════════
\section{Ưu \& Nhược Điểm}
% ══════════════════════════════════════════════════════════════════════════════

\subsubsection*{Ưu điểm}

\begin{itemize}
  \item \textbf{Pipeline đa tầng toàn diện:} Kết hợp cả 4 góc nhìn (RFM phân đoạn hành vi, ML dự
    đoán xác suất, KMeans phân nhóm, Kaplan-Meier đo tốc độ rời bỏ) --- mỗi bước bổ sung thông tin
    mà bước trước không có.
  \item \textbf{Dual preprocessing path:} Tách biệt rõ ràng tree vs linear pipeline, đảm bảo đúng
    methodology cho từng nhóm model.
  \item \textbf{Leakage-free KMeans:} KMeans chỉ fit trên pre-SMOTE training data, test set dùng
    \code{predict()} --- không có data leakage.
  \item \textbf{Risk Score thực tế:} Tích hợp \code{cashback\_pct} vào Risk Score phản ánh giá trị
    kinh tế của từng khách, không chỉ xác suất churn thuần túy.
  \item \textbf{Action Matrix có survival context:} Điều chỉnh ưu tiên theo Tenure
    (onboarding vs retention vs VIP) --- phù hợp hơn thực tế CSKH.
\end{itemize}

\subsubsection*{Nhược điểm}

\begin{itemize}
  \item \textbf{Tập dữ liệu nhỏ (5\,630 KH):} Silhouette score KMeans chỉ đạt 0.1183 (khá thấp)
    --- ranh giới giữa 5 clusters không sắc nét, dễ bị ảnh hưởng bởi outlier.
  \item \textbf{\code{cashback\_per\_order} là proxy không hoàn hảo:} Không có dữ liệu doanh thu
    hay LTV thực tế --- cashback có thể không phản ánh đúng giá trị kinh tế từng khách.
  \item \textbf{SMOTE tạo synthetic samples:} Trong tập nhỏ, SMOTE có thể tạo các điểm synthetic
    quá gần nhau, làm model overfit trên minority class.
  \item \textbf{RFM standalone yếu (AUC $= 0.55$):} Nếu dùng RFM độc lập để phân loại churn, hiệu
    quả rất thấp --- chỉ có giá trị khi kết hợp với ML.
  \item \textbf{Không có SHAP per-cluster:} SHAP chỉ tính global, chưa phân tích driver của từng
    cluster riêng biệt.
\end{itemize}

% ══════════════════════════════════════════════════════════════════════════════
\section{Hướng Phát Triển Tương Lai}
% ══════════════════════════════════════════════════════════════════════════════

\textbf{a. Đưa vào Production với MLflow / Model Registry}
\begin{itemize}
  \item Version control models, auto-retrain theo schedule (monthly/quarterly).
  \item A/B test giữa model versions trước khi deploy toàn phần.
\end{itemize}

\textbf{b. Feature Store \& Online Scoring}
\begin{itemize}
  \item Tách biệt batch scoring (monthly report) và real-time scoring (trigger khi có event: mua
    hàng mới / khảo sát mới).
  \item Pipeline inference $< 100$ms với serving framework (FastAPI + Redis cache).
\end{itemize}

\textbf{c. Explainability Dashboard}
\begin{itemize}
  \item Giao diện Streamlit/Dash hiển thị SHAP driver cho từng KH --- CSKH không cần đọc CSV.
  \item Filter theo Risk Tier, Cluster, khoảng thời gian.
\end{itemize}

\textbf{d. Causal Inference}
\begin{itemize}
  \item Thay thế ``correlation $\to$ action'' bằng ``uplift modeling'' --- đo lường xem can thiệp
    nào thực sự tạo ra sự khác biệt (treated vs control).
  \item Áp dụng Double ML hoặc causal forest để ước lượng CATE (Conditional Average Treatment
    Effect).
\end{itemize}

\textbf{e. Thu thập thêm dữ liệu}
\begin{itemize}
  \item Dữ liệu click-stream, session time, cart abandonment --- enrich feature space đáng kể.
  \item Dữ liệu revenue thực tế thay vì cashback proxy $\to$ Risk Score chính xác hơn.
\end{itemize}

\textbf{f. Cải thiện KMeans}
\begin{itemize}
  \item Thử Gaussian Mixture Model (GMM) để xử lý cluster overlap tốt hơn (silhouette hiện tại chỉ
    0.1352).
  \item Thêm SHAP per-cluster analysis để hiểu driver churn riêng của từng nhóm khách hàng.
\end{itemize}

\textbf{g. Time-series Features}
\begin{itemize}
  \item Hiện tại chỉ có snapshot feature. Tạo thêm rolling features (trend 3 tháng, seasonal
    pattern) nếu có transaction history.
\end{itemize}

\textbf{h. Semi-supervised Learning}
\begin{itemize}
  \item Dữ liệu có label churn thực sự khá ít (948 KH churn trong 5\,630). Thử semi-supervised
    approach để tận dụng unlabeled data từ các kỳ khác.
\end{itemize}

% ══════════════════════════════════════════════════════════════════════════════
\newpage
\appendix
\section{Phụ Lục --- Artifacts}
% ══════════════════════════════════════════════════════════════════════════════

\begin{center}
\begin{tabularx}{\textwidth}{llX}
  \toprule
  \textbf{Loại} & \textbf{File} & \textbf{Mô tả} \\
  \midrule
  Model & \code{models/E\_Commer\_Data/xgb\_clust\_model.pkl} & Production model \\
  Model & \code{models/E\_Commer\_Data/kmeans\_fe\_model.pkl} & KMeans $K=5$ \\
  Model & \code{models/E\_Commer\_Data/risk\_score\_meta.pkl} & Quintile thresholds \\
  Data  & \code{data/E\_Commer\_Data/model\_comparison.csv}  & 6-model metrics \\
  Data  & \code{data/E\_Commer\_Data/cluster\_summary.csv}   & 5 personas \\
  Data  & \code{data/E\_Commer\_Data/targeted\_action\_list.csv} & Action list đầy đủ \\
  \bottomrule
\end{tabularx}
\end{center}

\vfill
\begin{center}
  \small\itshape
  Báo cáo tổng hợp --- E-Commerce Churn Alert Pipeline\\
  XGB-Clust ROC-AUC = 0.974 $\mid$ $K=5$ $\mid$ Log-rank $p \approx 8.69 \times 10^{-9}$
\end{center}

\end{document}
